\section{Original homogeneous global sections are genuine fixed-model graded sections}
\label{sec:projective-native-homogeneous-model-sections}

\begin{theorem}[Original homogeneous coefficient and actual native generic germ {\cite[\S1.2]{BJ22}}]
\label{thm:projective-native-homogeneous-coefficient}
\lean{EnvelopeProperty.ProjectiveNativeHomogeneousModelSections.germ_homogeneousCoefficient_eq_aeval}
\leanok
\uses{thm:projective-native-model-global-section-values,lem:fc-original-homogeneous-section-coordinate,lem:original-native-coordinate-coefficients}
For any original nonempty coordinate chart, the native generic germ of the
whole original pulled homogeneous coefficient is the evaluation of its
original polynomial at the actual original native generic coordinate ratios.
\end{theorem}
\begin{proof}
The original whole projective chart section equivalence preserves the
original coefficient structure map, so its composition with the actual
native generic pullback is an original-field algebra map. Composing with
actual coordinate dehomogenization sends every original variable to the
actual original generic coordinate ratio. Polynomial algebra extensionality
identifies this map with evaluation at those ratios. The genuine original
homogeneous coefficient is exactly the dehomogenized polynomial section.
\end{proof}

\begin{theorem}[Original homogeneous forms give genuine whole model graded pieces {\cite[\S1.2]{BJ22}}]
\label{thm:projective-native-homogeneous-image}
\lean{EnvelopeProperty.ProjectiveNativeHomogeneousModelSections.homogeneousImage, EnvelopeProperty.ProjectiveNativeHomogeneousModelSections.homogeneousCoefficient, EnvelopeProperty.ProjectiveNativeHomogeneousModelSections.homogeneousImage_eq_monomial, EnvelopeProperty.ProjectiveNativeHomogeneousModelSections.homogeneousCoefficient_mem_piece, EnvelopeProperty.ProjectiveNativeHomogeneousModelSections.homogeneousCoefficient_chart_fraction}
\leanok
\uses{thm:projective-native-homogeneous-coefficient,lem:fc-native-rank-one-homogeneous-polynomial,lem:gl-sections}
Map every original variable to its literal original coordinate element in the
one fixed reconstructed homogeneous algebra. An original form F of degree d
maps to the actual monomial $g T^d$, where g is F evaluated at the original
pivot ratios in the entire native function field. Thus g belongs to the
actual fixed-model graded piece and hence its actual global section space.
In every nonempty original chart, its fraction $g/D_j$ is exactly the whole
original pulled homogeneous coefficient under the genuine chart equivalence.
\end{theorem}
\begin{proof}
Polynomial substitution sends actual degree-one coordinates to actual
degree-one monomials, so the image is homogeneous of degree d. On a rank-one
polynomial ring, every homogeneous polynomial of degree d is its unique
monomial of that degree. Evaluate the auxiliary variable at one: the original
substitution becomes evaluation of F at the actual original generic ratios,
which identifies its sole coefficient as g. Membership in the actual
homogeneous subalgebra proves membership of g in its graded piece. Apply
actual model dehomogenization. Its value on $g T^d$ is $g/D_j$, while its value on
each original coordinate is the original coordinate ratio. Genuine original
generic coordinate transition identities identify this with evaluation at
the j-chart native ratios. The preceding theorem identifies it with the
actual original coefficient germ, and the whole native chart equivalence
retains that germ literally. No section-realization or coefficient-matching
premise is supplied.
\end{proof}
